\documentclass[10pt,a4paper]{article}
\usepackage[margin=1.8cm]{geometry}
\usepackage[T1]{fontenc}
\usepackage[utf8]{inputenc}
\usepackage{enumitem,booktabs,array,xcolor}
\setlist{nosep,leftmargin=1.2em}
\pagestyle{empty}
\setlength{\parindent}{0pt}\setlength{\parskip}{4pt}
\begin{document}
{\Large\bfseries Spectral line diagnostics -- work split}\hfill{\small M2 IRT, MC ``Radiation'' 2026--27}

\textbf{Goal.} Tool that starts from a real spectrum, (1) fits the continuum, (2) finds and identifies the lines, (3) characterises them and derives the physical conditions of the gas. \textbf{Check-ins:} Oct 14, Nov 18. \textbf{Oral:} Dec 9 (50\% of grade, may be individual).

\textbf{Scope (updated after check-in 1, Oct 14).} Optical emission-line spectra of \emph{planetary nebulae}. Main target: NGC 7662 (HST/STIS G430L + G750L, MAST programme 12600, slit on the nebula); second target: NGC 7027 (programme 8285, richest red spectrum). The absorption work done before check-in 1 (Voigt profiles, equivalent widths, curves of growth, $(N,b)$ fit, all validated on synthetic spectra) is kept as a separate module.

\begin{tabular}{@{}p{8.2cm}@{\hspace{0.4cm}}p{8.2cm}@{}}
\toprule
\textbf{A -- data $\to$ lines} & \textbf{B -- lines $\to$ physics}\\
\midrule
\begin{itemize}
\item MAST download; 1D extraction from the 2D STIS image (sum the nebular rows, propagate errors)
\item Continuum fit and subtraction (sigma-clipping that rejects emission lines)
\item Nebular line list (H, He recombination lines; forbidden lines of O, N, S, Ne, Ar; air vs vacuum wavelengths)
\item Line detection and identification against the line list (radial velocity from the matched lines)
\item Interactive visualisation
\end{itemize}
&
\begin{itemize}
\item Gaussian line profiles including the instrumental resolution; multi-line fits of blends (H$\alpha$ + [N\,\textsc{ii}], H$\gamma$ + [O\,\textsc{iii}]\,4363)
\item Line fluxes with error propagation (noise and continuum)
\item Reddening from the Balmer decrement; dereddened fluxes
\item Electron temperature ([O\,\textsc{iii}], [N\,\textsc{ii}]), electron density ([S\,\textsc{ii}], [O\,\textsc{ii}]), ionic abundances
\item Uncertainties (Monte Carlo)
\end{itemize}\\
\bottomrule
\end{tabular}

\textbf{Interface.} A outputs a \texttt{Spectrum} (flux, error, continuum) plus a line table: name, $\lambda_\text{rest}$, observed centre guess, fitting window. B only needs this and returns a table of line fluxes and the derived physical quantities.

\textbf{Avoiding the dependency.} B first works on \emph{synthetic emission spectra} (known line fluxes on a known continuum) and on a hand-made line table for NGC 7662. A's automatic identification later produces the same table; B's code does not change.

\textbf{Shared work.} Repo and interfaces; synthetic spectrum generator; validation (inject known fluxes, check recovery; compare our temperatures and densities with PyNeb and with published values for the same nebulae); report and slides (each presents own half).

\textbf{Timeline.}\par
\begin{tabular}{@{}p{2.5cm}p{14.2cm}@{}}
\toprule
Until Oct 14 & Repo, interfaces, synthetic generator; absorption chain (profiles, EW, COG) validated on synthetic spectra. MAST search for planetary nebula spectra.\\
Oct 14 -- Nov 18 & A: extraction of NGC 7662, continuum, line list, detection and identification. B: Gaussian fits, fluxes, reddening, electron temperature. Demo on the real spectrum.\\
Nov 18 -- Dec 8 & Electron density, abundances, uncertainties, NGC 7027 as second object, interactive viz, comparison with the literature, slides and rehearsal.\\
Dec 9 & Oral presentation.\\
\bottomrule
\end{tabular}

\textbf{Working habits.} Shared git repo, personal branches merged into \texttt{main} by reviewed pull requests; 15--20 min weekly sync; integrate A's output into B's fits every 1--2 weeks; review each other's code so both understand the full pipeline.

\textbf{Status (see README progress log for details).} Done: \texttt{core} data classes, \texttt{profiles} (Gaussian and Voigt optical depth), \texttt{ew}, \texttt{io.make\_synthetic}, \texttt{linelist} (CIV, SiII, FeII), \texttt{cog}, with validation notebooks 00--04. Candidate planetary nebula spectra found in MAST; the pipeline 1D products are too noisy for an extended nebula, so we extract our own. Next: \texttt{io} reader and extraction for STIS \texttt{\_sx2} files, then \texttt{continuum}.
\end{document}
